\chapter{Analiza wymagań}
\label{chap:algs}

Analiza wymagań stanowi kluczowy etap w procesie projektowania i implementacji aplikacji UniGuesser. Celem tego rozdziału jest szczegółowe określenie funkcjonalności, jakie powinna spełniać aplikacja, a także zdefiniowanie kryteriów akceptacji, które pozwolą ocenić jej gotowość do wdrożenia i użytkowania. W ramach analizy wymagań uwzględniono zarówno potrzeby użytkowników końcowych, jak i aspekty techniczne związane z wyborem platformy oraz technologii implementacyjnych.

\section{Koncepcja gry (Karina Wołoszyn) }

UniGuesser to interaktywna aplikacja edukacyjno-rozrywkowa, której celem jest odgadywanie lokalizacji związanych z kampusem Politechniki Gdańskiej. Gra łączy elementy zabawy, edukacji oraz eksploracji przestrzeni uczelni, a dzięki technologii PWA użytkownicy mogą korzystać z niej na komputerach i urządzeniach mobilnych bez konieczności instalacji.

\bigskip
\noindent
\textbf{Celem UniGuesser jest:}
\begin{itemize}
    \item angażowanie graczy w interaktywną zabawę polegającą na identyfikacji lokalizacji na podstawie zdjęć,
    \item rozwijanie orientacji przestrzennej w obrębie kampusu uczelni,
    \item aktywne poznawanie nowych miejsc i życia akademickiego w formie przyjemnej i edukacyjnej.
\end{itemize}

\bigskip
\noindent
\textbf{Aplikacja jest skierowana do:}
\begin{itemize}
    \item studentów Politechniki Gdańskiej, w tym nowych studentów poznających kampus,
    \item pracowników uczelni,
    \item osób odwiedzających kampus, np. podczas dni otwartych,
    \item entuzjastów gier geolokalizacyjnych i edukacyjnych, którzy cenią naukę przez zabawę.
\end{itemize}

\bigskip
\noindent
\textbf{UniGuesser może pełnić funkcję:}
\begin{itemize}
    \item narzędzia integracyjnego dla studentów i pracowników uczelni,
    \item elementu wydarzeń akademickich, takich jak dni otwarte czy spotkania organizacyjne,
    \item gry terenowej zachęcającej do aktywnego zwiedzania kampusu.
\end{itemize}

\bigskip
\noindent
\textbf{Zaletami aplikacji UniGuesser są między innymi:}
\begin{itemize}
    \item dostępność na różnych urządzeniach dzięki technologii PWA,
    \item połączenie zabawy z nauką i poznawaniem kampusu,
    \item motywacja do rywalizacji poprzez system punktacji i rankingów,
    \item zachęta do aktywnego odkrywania otoczenia w trybie terenowym,
    \item wsparcie integracji społeczności akademickiej i współzawodnictwa między graczami.
\end{itemize}

\section{Przebieg gry (Karina Wołoszyn)}

Rozgrywka w UniGuesser przebiega w kilku etapach, które zapewniają płynność gry i umożliwiają graczowi śledzenie swoich wyników:

\begin{enumerate}
    \item \textbf{Rozpoczęcie gry} - użytkownik wybiera tryb i trudność rozgrwyki, a w przypadku gdy nie jest zalogowany - wpisuje swój nick. System umożliwia także kontynuację poprzedniej gry wraz z zapisanymi punktami i historią odpowiedzi lub rozpoczęcie nowej sesji.
    \item \textbf{Losowanie lokacji} - aplikacja losowo wybiera zdjęcie lokalizacji z bazy danych wraz z jej współrzędnymi, przy czym w danej sesji nie powtarzają się już wyświetlone miejsca.
    \item \textbf{Odgadywanie lokalizacji} - gracz wskazuje na interaktywnej mapie miejsce odpowiadające wylosowanemu zdjęciu. Interfejs pozwala przybliżać i oddalać mapę oraz przesuwać ją w celu precyzyjnego wyboru. Gracz może wielokrotnie zaznaczać punkty, jednak decyduje ostatnie kliknięcie.
    \item \textbf{Zatwierdzenie wyboru} - gracz potwierdza swój wybór, klikając przycisk zatwierdzający decyzję.
    \item \textbf{Wyświetlenie wyniku} - po zatwierdzeniu odpowiedzi system oblicza odległość między wskazanym a faktycznym miejscem, przyznaje odpowiednią liczbę punktów i pokazuje prawidłowe położenie lokalizacji na mapie.
    \item \textbf{Podsumowanie rundy} - gracz otrzymuje informacje o zdobytych punktach w danej rundzie oraz może przeglądać krótką statystykę swoich postępów.
    \item \textbf{Kontynuacja lub zakończenie gry} - jeśli gracz nie rozegrał jeszcze wszystkich rund, system pozwala przejść do kolejnej. W przeciwnym wypadku wyświetlane jest podsumowanie całej gry w formie przejrzystej tabeli z możliwością podglądu wyników każdej rundy.
\end{enumerate}

Taka struktura rozgrywki sprawia, że gra jest czytelna, intuicyjna i pozwala zarówno na zabawę, jak i naukę poprzez aktywne poznawanie kampusu Politechniki Gdańskiej.

\section{Tryby rozgrywki (Karina Wołoszyn)}

Gra oferuje dwa tryby rozgrywki, które urozmaicają zgadywanie lokalizacji. Oprócz trybu podstawowego, rozgrywanego w miejscu, dostępny jest również tryb terenowy, przeznaczony do gry na zewnątrz i w ruchu. Poniżej przedstawiono szczegółowy opis trybu terenowego.

Tryb terenowy pobiera lokalizację gracza za pomocą sygnału GPS dostępnego w urządzeniu przenośnym i wyświetla ją na mapie. Gracz może w prosty sposób zmieniać widok na ekranie, przełączając się między mapą a zdjęciem. Tryb terenowy posiada również system podpowiedzi oparty na mechanice „ciepło-zimno”, mający na celu ograniczenie potencjalnej frustracji gracza. Podpowiedź może być użyta maksymalnie trzy razy i działa poprzez wyświetlanie odpowiednich kolorów: czerwonego, jeśli gracz zbliża się do celu w ciągu ostatnich 30 sekund, lub niebieskiego, jeśli się oddala. Gdy gracz znajduje się w promieniu 5 metrów od miejsca docelowego i zatwierdzi swój „strzał”, system kończy rundę i wyświetla wynik tak jak w trybie podstawowym.

W trybie terenowym kluczowa jest dokładność użytkownika. Poprzez aktywne poruszanie się i poszukiwanie miejsca na podstawie jego widoku, gracz w praktyczny sposób poznaje teren kampusu politechniki oraz lokalizacje związane z uczelnią.


\section{Wymagania funkcjonalne (Karina Wołoszyn)}

Wymagania funkcjonalne określają kluczowe funkcje, które aplikacja musi realizować. Definiują one, jak system powinien się zachowywać w określonych sytuacjach, jakie działania ma wykonywać oraz jakie rezultaty ma osiągać w odpowiedzi na określone dane wejściowe.

Fundamentalnym wymogiem całego projektu było sprecyzowanie wymagań dotyczących podstawowej rozgrywki. To właśnie w tym trybie gracz spędza najwięcej czasu, dlatego kluczowe jest, aby była ona dopracowana, angażująca i zachęcająca do ponownego grania. W związku z tym wyróżniono następujące wymagania:

\begin{itemize}
    \item \textbf{Adekwatna punktacja} - aby gracz czuł się odpowiednio nagrodzony lub ukarany za swoją odpowiedź, system powinien przydzielać punkty proporcjonalnie do odległości pomiędzy faktycznym miejscem a wskazaniem gracza. Aplikacja musi zatem nie tylko obliczać tę odległość, lecz także przeliczać ją na odpowiednią liczbę punktów.
    \item \textbf{Zapisywanie wyników pojedynczego gracza} - wyniki każdej rozgrywki są zapisywane, aby gracz mógł w przyszłości przeglądać swoje wcześniejsze sesje. Podgląd wyników umożliwia analizę i rozwój umiejętności, a także pozwala zobaczyć, gdzie dokładnie padły odpowiedzi na mapie wraz z informacjami o lokalizacjach.
    \item \textbf{Tryby gry} - rozgrywka posiada dwa tryby gry: klasyczny (losowe lokacje z bazy danych) oraz tematyczny (lokacje powiązane z wybranym obszarem, np. kampusem uczelni lub konkretnym miastem).
    \item \textbf{Rejestracja i logowanie} - system powinien rozróżniać role użytkowników poprzez proces rejestracji i logowania. Zwykły użytkownik ma ograniczone uprawnienia w porównaniu z Moderatorem czy Administratorem. Nowe konto można założyć za pomocą prostego formularza rejestracyjnego.
    \item \textbf{Dodawanie nowych lokacji przez użytkowników} - zarejestrowani użytkownicy mogą zgłaszać nowe lokacje do bazy danych, dodając opis, współrzędne geograficzne oraz zdjęcie. Propozycje trafiają do weryfikacji przed publikacją.
    \item \textbf{Nadzór lokacji} - role takie jak Administrator czy Moderator mają możliwość zatwierdzania, edytowania lub odrzucania dodanych lokalizacji, aby zapobiec nadużyciom, naruszeniom praw autorskich oraz publikacji niepożądanych treści.
    \item \textbf{Ranking graczy} - ranking zwiększa element rywalizacji między graczami i wspiera rozwój społeczności. Przedstawia on nick gracza, jego wynik, tryb rozgrywki oraz poziom trudności. Ranking posiada paginację i ograniczenie do jednego rekordu dla danego zestawu ustawień, co usprawnia jego działanie.
    \item \textbf{Prosty i intuicyjny interfejs} - użytkownik powinien już podczas pierwszego kontaktu z systemem intuicyjnie rozumieć jego działanie. Projekt interfejsu opiera się na minimalizmie i spójności wizualnej w całej aplikacji.
    \item \textbf{Kolorystyka interfejsu} - zastosowana kolorystyka powinna być zgodna z barwami Politechniki Gdańskiej, zachowując estetykę i czytelność.
    \item \textbf{Lokacje związane z życiem uczelnianym} - baza danych powinna zawierać lokacje związane z Politechniką Gdańską i życiem studenckim, np. budynki uczelni, akademiki, kawiarnie, parki czy znane miejsca spotkań studentów.
    \item \textbf{Responsywne UI} - interfejs powinien poprawnie działać na komputerach oraz urządzeniach mobilnych, dostosowując układ elementów do rozmiaru ekranu, aby zachować komfort gry niezależnie od urządzenia.
    \item \textbf{Responsywna mapa} - mapa powinna umożliwiać graczowi umieszczanie pinezek jako odpowiedzi, a także przybliżanie, oddalanie, przesuwanie i automatyczne wyśrodkowanie na docelową lokację po zakończeniu rundy.
    \item \textbf{Aplikacja PWA} - aplikacja powinna wspierać instalację w formie PWA, umożliwiając korzystanie z niej offline i z poziomu ekranu głównego urządzenia.
    \item \textbf{Kontynuacja gry} - system powinien pozwalać użytkownikowi na kontynuację niedokończonej rozgrywki z zachowaniem dotychczasowych punktów oraz historii odpowiedzi.
    \item \textbf{Hamburger menu} - przy małym rozmiarze ekranu opcje z górnego paska menu powinny być zwijane do tzw. hamburger menu, przy zachowaniu kolejności i logiki nawigacji.
    \item \textbf{System podpowiedzi „ciepło-zimno”} - w trybie terenowym gracz ma dostęp do systemu podpowiedzi, który pomaga mu zlokalizować docelowe miejsce, nie ujawniając go wprost.
\end{itemize}

Wymagania te stanowią podstawę do stworzenia aplikacji spełniającej główne założenia projektu, a jednocześnie podnoszącej jakość doświadczenia i satysfakcję użytkownika.


\section{Wymagania niefunkcjonalne (Filip Jezierski)}
Wymagania niefunkcjonalne określają kryteria jakościowe, które aplikacja powinna spełniać, aby zapewnić odpowiednie doświadczenie użytkownika oraz efektywne działanie systemu. W kontekście opracowywanej aplikacji mobilnej, kluczowe wymagania niefunkcjonalne obejmują aspekty takie jak wydajność, bezpieczeństwo oraz użyteczność.

\bigskip
\noindent
\textbf{Wymagania jakościowe}
\begin{itemize}
    \item obsługa różnych rozdzielczości ekranów i urządzeń mobilnych,
    \item dostępność aplikacji na popularnych platformach mobilnych (Android, iOS), a także jako aplikacja webowa,
    \item intuicyjny i responsywny interfejs użytkownika,
    \item aplikacja powinna w odpowiedni sposób obsługiwać błędy i wyjątki, takie jak nieprawidłowe dane wejściowe czy brak internetu, zapewniając stabilność działania,
    \item ochrona danych użytkowników poprzez implementację odpowiednich mechanizmów bezpieczeństwa - szyfrowanie danych, bezpieczne logowanie, ochrona przed atakami typu SQL Injection i XSS.
\end{itemize}

\bigskip
\noindent
\textbf{Wymagania wydajnościowe}
\begin{itemize}
    \item szybkie (<2 s) ładowanie aplikacji i minimalny czas oczekiwania na reakcję interfejsu użytkownika,
    \item ładowanie zdjęć lokalizacji w czasie nie dłuższym niż 3 sekundy,
    \item szybkie (<10 s) ładowanie danych geolokalizacyjnych zapewniające płynne działanie trybu gry opartego na lokalizacji,
\end{itemize}

\section{Wybór platformy (Filip Jezierski)}
Wybór odpowiedniej platformy dla opracowywanej aplikacji mobilnej jest kluczowym elementem procesu projektowego, mającym istotny wpływ na jej funkcjonalność, dostępność oraz doświadczenie użytkownika, a także łatwość wdrożenia i utrzymania. Na etapie analizy wymagań podjęto decyzję o stworzeniu dwóch testowych wersji aplikacji: aplikacji webowej typu PWA (Progressive Web App) oraz natywnej aplikacji mobilnej wykorzystującej technologię .NET MAUI (Multi-platform App UI). Obie platformy oferują unikalne korzyści i wyzwania, które zostały dokładnie rozważone w kontekście specyfiki projektu.

\subsection{Natywna aplikacja mobilna}
Aplikacje natywne są tworzone z wykorzystaniem dedykowanych narzędzi i języków programowania przeznaczonych dla konkretnych systemów operacyjnych, takich jak Android czy iOS~\cite{dotnet-maui-docs}. W przypadku technologii .NET MAUI możliwe jest tworzenie aplikacji wieloplatformowych przy zachowaniu wspólnej bazy kodu, co znacząco skraca czas implementacji i ułatwia utrzymanie.

Jednakże proces publikacji aplikacji natywnych w sklepach takich jak Google Play czy Apple App Store wiąże się z dodatkowymi wymaganiami i procedurami zatwierdzania, a także z koniecznością budowania i testowania aplikacji na różne systemy operacyjne, co może wydłużyć czas wdrożenia.

\subsection{Aplikacja webowa PWA}
Aplikacja webowa typu PWA (Progressive Web App) łączy cechy klasycznej strony internetowej z funkcjonalnością aplikacji mobilnej~\cite{pwa-docs}. Dzięki wykorzystaniu nowoczesnych technologii webowych, możliwe jest stworzenie aplikacji działającej w trybie offline, instalowanej bezpośrednio z przeglądarki oraz zapewniającej wrażenia zbliżone do aplikacji natywnej.

Ze względu na łatwość dostępu (poprzez przeglądarkę internetową) oraz brak konieczności przechodzenia przez proces zatwierdzania w sklepach aplikacji, PWA stanowi atrakcyjną opcję dla naszego projektu. Jednakże, pewne ograniczenia związane z dostępem do funkcji urządzenia (takich jak geolokalizacja) oraz wydajnością w porównaniu do aplikacji natywnych muszą być uwzględnione podczas projektowania i implementacji.

\section{Ograniczenia projektu (Filip Jezierski)}
Ograniczenia projektu wynikają przede wszystkim z przyjętych założeń technologicznych oraz charakteru opracowywanego rozwiązania. Kluczowe znaczenie mają tu specyfika aplikacji webowej typu PWA, zastosowanie mechanizmów geolokalizacji jako podstawowego elementu funkcjonalnego, a także decyzje dotyczące sposobu pozyskiwania i wykorzystania materiałów wizualnych. Wymienione czynniki w naturalny sposób wpływają na zakres możliwości implementacyjnych, wydajność rozwiązania oraz doświadczenie użytkownika końcowego.

\subsection{Środowisko aplikacji}
Aplikacja będzie udostępniana jako aplikacja webowa, dostępna przez przeglądarkę internetową, z możliwością instalacji na ekranie głównym urządzenia. Funkcjonalność trybu gry wykorzystującego geolokalizację będzie dostępna wyłącznie na urządzeniach mobilnych, które wspierają tę technologię. W przypadku urządzeń stacjonarnych, nawet jeśli geolokalizacja jest technicznie możliwa, jej dokładność jest niewystarczająca do zapewnienia poprawnego działania tej funkcji.

\subsection{Zdjęcia lokalizacji}
W projekcie zrezygnowano z wykorzystywania zdjęć 360° pochodzących z zewnętrznych źródeł, takich jak interfejsy API udostępniane przez firmy trzecie. Decyzję tę podjęto z uwagi na chęć uniknięcia potencjalnych problemów związanych z licencjonowaniem, ograniczoną dostępnością materiałów oraz zróżnicowaną jakością zdjęć. Dodatkowo korzystanie z takich źródeł mogłoby wprowadzać ograniczenia dotyczące liczby i różnorodności dostępnych lokalizacji.

W celu uproszczenia procesu pozyskiwania materiałów wizualnych oraz umożliwienia użytkownikom współtworzenia bazy lokalizacji, aplikacja będzie wykorzystywać standardowe zdjęcia dwuwymiarowe. Tego typu fotografie nie wymagają specjalistycznego sprzętu i mogą być wykonywane przy użyciu dowolnego smartfona. Użytkownicy będą mieli możliwość przesyłania własnych zdjęć lokalizacji, które po weryfikacji przez moderatorów zostaną dodane do bazy danych aplikacji.

Takie rozwiązanie zapewnia pełną kontrolę nad jakością i zgodnością materiałów z założeniami projektu, a jednocześnie pozwala na stopniowe rozszerzanie bazy lokalizacji o nowe zdjęcia, dostarczane zarówno przez zespół projektowy, jak i społeczność użytkowników.

\section{Kryteria akceptacji (Filip Jezierski)}
Kryteria akceptacji stanowią zestaw wymagań, które muszą zostać spełnione, aby aplikacja mogła zostać uznana za gotową do wdrożenia i użytkowania. Kryteria te obejmują zarówno aspekty funkcjonalne, jak i niefunkcjonalne, zapewniając kompleksową ocenę jakości i użyteczności opracowanego rozwiązania.

\bigskip
\noindent
\textbf{Ogólne kryteria akceptacji:}
\begin{itemize}
    \item Aplikacja działa na urządzeniach mobilnych z systemami Android i iOS, na komputerach stacjonarnych oraz w przeglądarkach internetowych
    \item Intuicyjny interfejs użytkownika umożliwiający łatwe korzystanie z aplikacji
    \item Responsywny design dostosowujący się do różnych rozmiarów ekranów, w szczególności urządzeń mobilnych
    \item Stabilne działanie aplikacji bez krytycznych błędów i awarii
    \item Testy manualne wszystkich zaplanowanych funkcjonalności przeprowadzone na wszystkich wspieranych platformach
\end{itemize}

\bigskip
\noindent
\textbf{Zaimplementowane funkcjonalności:}
\begin{itemize}
    \item Rejestracja i logowanie użytkowników
    \item Wybór trybu rozgrywki i poziomu trudności
    \item Dwa tryby gry:
          \begin{itemize}
              \item odgadywanie lokalizacji na podstawie zdjęcia poprzez wskazanie jej na mapie
              \item odnajdywanie lokalizacji na podstawie zdjęcia i potwierdzanie obecności poprzez geolokalizację
          \end{itemize}
    \item Możliwość przerywania i kontynuowania rozgrywki
    \item Wyświetlanie podsumowania wyników każdej rundy gry
    \item Historia rozgrywek dostępna w profilu użytkownika
    \item Ogólnodostępny ranking graczy
    \item Dodawanie lokalizacji przez użytkowników (po weryfikacji przez moderatorów)
\end{itemize}
